%%%%%%%%%%%%%%%%%%%%%%%%%%%%%%%%%%%%%%%%%%%%%%%%%%%%%%%%%%%%%%%%%%%%%%%%%%%%%%
% Guide de référence CarnoTeX
%%%%%%%%%%%%%%%%%%%%%%%%%%%%%%%%%%%%%%%%%%%%%%%%%%%%%%%%%%%%%%%%%%%%%%%%%%%%%%

\documentclass[a4paper,11pt]{article}

\usepackage[all]{carnotex}
\level{notes}

%=================================================
% Commandes propres au guide
%=================================================

\newcommand{\pkg}[1]{\texttt{#1}}
\newcommand{\cmd}[1]{\texttt{\char`\\#1}}
\newcommand{\argu}[1]{\texttt{\char`\{#1\char`\}}}
\newcommand{\opt}[1]{\texttt{[#1]}}

\NewDocumentCommand{\syntaxe}{+m}{%
  \par\noindent
  \texttt{\detokenize{#1}}%
  \par
}

\NewDocumentEnvironment{commande}{m m +b}{%
  \subsection{\cmd{#1}}
  \cadre{CarnoGrey}[Package]{\pkg{#2}}
  #3
}{}

\NewDocumentCommand{\Description}{+m}{%
  \cadre{CarnoDef}[Description]{#1}
}

\NewDocumentCommand{\Syntaxe}{+m}{%
  \cadre{CarnoRem}[Syntaxe]{#1}
}

\NewDocumentCommand{\Arguments}{+m}{%
  \cadre{CarnoProp}[Arguments et options]{#1}
}

\NewDocumentCommand{\Usage}{+m}{%
  \cadre{CarnoExe}[Philosophie d'usage]{#1}
}

\NewDocumentCommand{\Exemple}{+m}{%
  \cadre{CarnoGreen}[Exemple de rendu]{#1}
}

\NewDocumentCommand{\Remarques}{+m}{%
  \cadre{CarnoYellow}[Remarques]{#1}
}

%=================================================
% Document
%=================================================

\title{Guide de référence\\\Large CarnoTeX}
\author{Yann Grizonnet}
\date{\today}

\begin{document}

\maketitle
\tableofcontents
\newpage

%%%%%%%%%%%%%%%%%%%%%%%%%%%%%%%%%%%%%%%%%%%%%%%%%%%%%%%%%%%%%%%%%%%%%%%%%%%%%%
\section{Présentation générale}
%%%%%%%%%%%%%%%%%%%%%%%%%%%%%%%%%%%%%%%%%%%%%%%%%%%%%%%%%%%%%%%%%%%%%%%%%%%%%%

CarnoTeX est un ensemble de packages LaTeX personnels destinés à produire des documents pédagogiques de mathématiques cohérents, lisibles à l'écran, projetables et imprimables.

\cadre{CarnoDef}[Architecture générale]{
\begin{itemize}
  \item \pkg{carnotex.sty} : socle graphique et technique ;
  \item \pkg{carnomath.sty} : commandes mathématiques ;
  \item \pkg{carnopeda.sty} : outils pédagogiques ;
  \item \pkg{carnolegacy.sty} : compatibilité avec les anciens documents.
\end{itemize}
}

\cadre{CarnoProp}[Options de chargement]{
\begin{itemize}
  \item \cmd{usepackage}\argu{carnotex} : noyau seul ;
  \item \cmd{usepackage}\opt{math}\argu{carnotex} : noyau + mathématiques ;
  \item \cmd{usepackage}\opt{peda}\argu{carnotex} : noyau + pédagogie ;
  \item \cmd{usepackage}\opt{legacy}\argu{carnotex} : noyau + compatibilité ;
  \item \cmd{usepackage}\opt{all}\argu{carnotex} : tout charger.
\end{itemize}
}

\cadre{CarnoRem}[Principes de conception]{
\begin{itemize}
  \item Une commande publique doit avoir une responsabilité claire.
  \item Les commandes internes sont préfixées par \cmd{Carno@}.
  \item Les anciennes commandes sont conservées uniquement dans \pkg{carnolegacy}.
  \item Toute nouvelle commande publique doit être testée et documentée.
\end{itemize}
}

%%%%%%%%%%%%%%%%%%%%%%%%%%%%%%%%%%%%%%%%%%%%%%%%%%%%%%%%%%%%%%%%%%%%%%%%%%%%%%
\section{Palette graphique}
%%%%%%%%%%%%%%%%%%%%%%%%%%%%%%%%%%%%%%%%%%%%%%%%%%%%%%%%%%%%%%%%%%%%%%%%%%%%%%

\subsection{Palette de base}

\noindent
\begin{tikzpicture}
  \foreach \couleur/\i in {
    CarnoBlue/0,
    CarnoCyan/1,
    CarnoGreen/2,
    CarnoYellow/3,
    CarnoOrange/4,
    CarnoRed/5,
    CarnoPurple/6,
    CarnoGrey/7
  }{
    \fill[\couleur] ({\i*\linewidth/8},0)
      rectangle ({(\i+1)*\linewidth/8},1.1);
    \node[anchor=north,font=\scriptsize\ttfamily]
      at ({(\i+0.5)*\linewidth/8},-0.15) {\couleur};
  }
\end{tikzpicture}

\bigskip

\subsection{Couleurs sémantiques}

\begin{center}
\renewcommand{\arraystretch}{1.4}
\begin{tabular}{|c|l|}
\hline
\textbf{Couleur} & \textbf{Usage} \\
\hline
\colorbox{CarnoDef}{\color{white}\rule{2cm}{0.2cm}} & \texttt{CarnoDef} : définitions \\
\hline
\colorbox{CarnoProp}{\color{white}\rule{2cm}{0.2cm}} & \texttt{CarnoProp} : propriétés \\
\hline
\colorbox{CarnoThm}{\color{white}\rule{2cm}{0.2cm}} & \texttt{CarnoThm} : théorèmes \\
\hline
\colorbox{CarnoExe}{\color{white}\rule{2cm}{0.2cm}} & \texttt{CarnoExe} : exemples \\
\hline
\colorbox{CarnoExo}{\color{white}\rule{2cm}{0.2cm}} & \texttt{CarnoExo} : exercices \\
\hline
\colorbox{CarnoPrb}{\color{white}\rule{2cm}{0.2cm}} & \texttt{CarnoPrb} : problèmes \\
\hline
\colorbox{CarnoRem}{\color{white}\rule{2cm}{0.2cm}} & \texttt{CarnoRem} : remarques \\
\hline
\colorbox{CarnoDem}{\color{white}\rule{2cm}{0.2cm}} & \texttt{CarnoDem} : démonstrations \\
\hline
\colorbox{CarnoSol}{\color{white}\rule{2cm}{0.2cm}} & \texttt{CarnoSol} : solutions \\
\hline
\end{tabular}
\end{center}

%%%%%%%%%%%%%%%%%%%%%%%%%%%%%%%%%%%%%%%%%%%%%%%%%%%%%%%%%%%%%%%%%%%%%%%%%%%%%%
\section{Mise en page}
%%%%%%%%%%%%%%%%%%%%%%%%%%%%%%%%%%%%%%%%%%%%%%%%%%%%%%%%%%%%%%%%%%%%%%%%%%%%%%

\begin{commande}{cadre}{carnotex}

\Description{
Affiche un cadre coloré pleine largeur, avec un titre optionnel placé au-dessus du cadre.
}

\Syntaxe{
\syntaxe{\cadre{couleur}{contenu}}
\syntaxe{\cadre{couleur}[titre]{contenu}}
}

\Arguments{
\begin{itemize}
  \item \texttt{couleur} : couleur du cadre ;
  \item \texttt{titre} : optionnel ;
  \item \texttt{contenu} : contenu du cadre.
\end{itemize}
}

\Exemple{
\cadre{CarnoDef}[Définition importante]{
Voici un cadre avec titre.
}

\cadre{CarnoProp}{
Voici un cadre sans titre.
}
}

\end{commande}

\begin{commande}{col}{carnotex}

\Description{
Place deux blocs côte à côte dans deux minipages.
}

\Syntaxe{
\syntaxe{\col{gauche}{droite}}
\syntaxe{\col[0.3][0.65]{gauche}{droite}}
}

\Arguments{
\begin{itemize}
  \item première option : largeur relative de la colonne gauche ;
  \item deuxième option : largeur relative de la colonne droite ;
  \item \texttt{gauche} : contenu de gauche ;
  \item \texttt{droite} : contenu de droite.
\end{itemize}
}

\Exemple{
\col{
Colonne de gauche.
}{
Colonne de droite.
}
}

\end{commande}

\begin{commande}{bordure}{carnotex}

\Description{
Active une bordure verticale colorée sur le bord droit des pages suivantes.
}

\Syntaxe{
\syntaxe{\bordure{titre}}
\syntaxe{\bordure[couleur]{titre}}
}

\Arguments{
\begin{itemize}
  \item \texttt{couleur} : couleur optionnelle ;
  \item \texttt{titre} : texte vertical affiché dans la bordure.
\end{itemize}
}

\Remarques{
La commande est persistante : une fois appelée, elle s'applique aux pages suivantes jusqu'à modification.
}

\end{commande}

%%%%%%%%%%%%%%%%%%%%%%%%%%%%%%%%%%%%%%%%%%%%%%%%%%%%%%%%%%%%%%%%%%%%%%%%%%%%%%
\section{Éléments graphiques}
%%%%%%%%%%%%%%%%%%%%%%%%%%%%%%%%%%%%%%%%%%%%%%%%%%%%%%%%%%%%%%%%%%%%%%%%%%%%%%

\begin{commande}{colorline}{carnotex}

\Description{
Trace une ligne horizontale.
}

\Syntaxe{
\syntaxe{\colorline}
\syntaxe{\colorline[couleur]}
\syntaxe{\colorline*[couleur][longueur][épaisseur]}
}

\Arguments{
\begin{itemize}
  \item étoile : centre le trait ;
  \item \texttt{couleur} : noire par défaut ;
  \item \texttt{longueur} : \cmd{linewidth} par défaut ;
  \item \texttt{épaisseur} : \texttt{0.5pt} par défaut.
\end{itemize}
}

\Exemple{
\colorline[CarnoProp]

\colorline*[CarnoThm][6cm][1pt]
}

\end{commande}

\begin{commande}{decoline}{carnotex}

\Description{
Trace une ligne décorative courte avec un motif central.
}

\Syntaxe{
\syntaxe{\decoline}
}

\Arguments{
Cette commande ne prend aucun argument.
}

\Exemple{
\decoline
}

\end{commande}

\begin{commande}{histo}{carnotex}

\Description{
Met en valeur un commentaire historique ou culturel.
}

\Syntaxe{
\syntaxe{\histo{contenu}}
}

\Arguments{
Un unique argument contenant le texte.
}

\Exemple{
\histo{
Petit commentaire historique.
}
}

\end{commande}

\begin{commande}{image}{carnotex}

\Description{
Insère une image depuis le dossier \texttt{img/} en conservant ses proportions.
}

\Syntaxe{
\syntaxe{\image{nom}}
\syntaxe{\image[options]{nom}}
\syntaxe{\image*{nom}}
\syntaxe{\image*[options]{nom}}
}

\Arguments{
\begin{itemize}
  \item \texttt{nom} : nom du fichier image dans le dossier \texttt{img/} ;
  \item \texttt{options} : options supplémentaires transmises à \cmd{includegraphics}.
\end{itemize}
}

\Usage{
La version non étoilée prend toute la largeur disponible.

La version étoilée utilise au maximum l'espace vertical restant sur la page.
}

\Exemple{
\IfFileExists{img/img2.png}{
   \image[width=4cm]{img2} \image[height=3cm]{img2}
   
   \image*{img2}
}{
  Exemple à compiler avec un fichier \texttt{img/img2.png}.
}
}

\end{commande}

\begin{commande}{carreaux}{carnotex}

\Description{
Dessine un quadrillage de type cahier avec marge rouge.
}

\Syntaxe{
\syntaxe{\carreaux{hauteur}{largeur}}
}

\Arguments{
\begin{itemize}
  \item \texttt{hauteur} : hauteur du quadrillage ;
  \item \texttt{largeur} : largeur du quadrillage.
\end{itemize}
}

\Exemple{
\carreaux{4}{6}
}

\end{commande}

%%%%%%%%%%%%%%%%%%%%%%%%%%%%%%%%%%%%%%%%%%%%%%%%%%%%%%%%%%%%%%%%%%%%%%%%%%%%%%
\section{Outils pratiques}
%%%%%%%%%%%%%%%%%%%%%%%%%%%%%%%%%%%%%%%%%%%%%%%%%%%%%%%%%%%%%%%%%%%%%%%%%%%%%%

\begin{commande}{rl}{carnotex}

\Description{
Ajoute une règle invisible pour augmenter la hauteur d'une ligne, notamment dans un tableau.
}

\Syntaxe{
\syntaxe{\rl}
}

\Arguments{
Aucun argument.
}

\Exemple{
\begin{tabular}{|c|c|}
\hline
$\dfrac{x^3}{3}$ & $\dfrac{x^3\rl}{3\rl}+k$ \rl\\\hline
$x^2$ & $x^3$ \\\hline
$x^2$ & $x^3$ \rl\\\hline
$\sqrt{x}$ & $\sqrt{x\rl}$\rl\\\hline
\end{tabular}
}

\end{commande}

\begin{commande}{restepage}{carnotex}

\Description{
Calcule l'espace vertical restant sur la page.
}

\Syntaxe{
\syntaxe{\restepage}
}

\Usage{
Utilisé notamment par \cmd{image*}.
}

\Exemple{
Espace restant : \the\restepage
}

\end{commande}

\begin{commande}{siecle}{carnotex}

\Description{
Compose proprement les siècles.
}

\Syntaxe{
\syntaxe{\siecle{21}}
\syntaxe{\siecle*{3}}
}

\Arguments{
\begin{itemize}
  \item argument obligatoire : numéro du siècle ;
  \item étoile : siècle avant Jésus-Christ.
\end{itemize}
}

\Exemple{
\siecle{21}

\siecle*{3}
}

\end{commande}

%%%%%%%%%%%%%%%%%%%%%%%%%%%%%%%%%%%%%%%%%%%%%%%%%%%%%%%%%%%%%%%%%%%%%%%%%%%%%%
\section{Mathématiques}
%%%%%%%%%%%%%%%%%%%%%%%%%%%%%%%%%%%%%%%%%%%%%%%%%%%%%%%%%%%%%%%%%%%%%%%%%%%%%%

\begin{commande}{N, Z, D, Q, R, Cp}{carnomath}

\Description{
Commandes courtes pour les ensembles usuels.
}

\Syntaxe{
\syntaxe{\[\N,\quad \Z,\quad \D,\quad \Q,\quad \R,\quad \Cp\]}
}

\Arguments{
Aucun argument.
}

\Exemple{
\[
\N,\quad \Z,\quad \D,\quad \Q,\quad \R,\quad \Cp
\]
}

\end{commande}

\begin{commande}{vc, coord, coor, vct}{carnomath}

\Description{
Commandes de confort pour les vecteurs et coordonnées.
}

\Syntaxe{
\syntaxe{\[\vc{AB}\qquad \coord{2}{3}\qquad \coor{1}{2}{3}\qquad \vct{u}{2}{-1}\]}
}

\Arguments{
\begin{itemize}
  \item \cmd{vc}\argu{nom} : vecteur ;
  \item \cmd{coord}\argu{x}\argu{y} : coordonnées en dimension 2 ;
  \item \cmd{coor}\argu{x}\argu{y}\argu{z} : coordonnées en dimension 3 ;
  \item \cmd{vct}\argu{nom}\argu{x}\argu{y} : vecteur avec coordonnées.
\end{itemize}
}

\Exemple{
\[
\vc{AB}
\qquad
\coord{2}{3}
\qquad
\coor{1}{2}{3}
\qquad
\vct{u}{2}{-1}
\]
}

\end{commande}

\begin{commande}{Oij, Oijk}{carnomath}

\Description{
Commandes pour les repères usuels du plan et de l'espace.
}

\Syntaxe{
\syntaxe{\Oij}
\syntaxe{\Oijk}
}

\Arguments{
Aucun argument.
}

\Exemple{
\Oij

\Oijk
}

\end{commande}

\begin{commande}{df, prop, thrm, lem, coro}{carnomath}

\Description{
Commandes de théorèmes numérotés avec version étoilée encadrée.
}

\Syntaxe{
\syntaxe{\df{contenu}}
\syntaxe{\df*{contenu}}
\syntaxe{\prop{contenu}}
\syntaxe{\prop*{contenu}}
\syntaxe{\thrm{contenu}}
\syntaxe{\thrm*{contenu}}
}

\Arguments{
Un argument obligatoire contenant l'énoncé.

La version étoilée utilise un cadre coloré.
}

\Exemple{
\df{
Une fonction affine est une fonction de la forme $f(x)=ax+b$.
}

\prop*{
La somme de deux fonctions continues est continue.
}

\thrm*{
Théorème de Pythagore.
}
}

\end{commande}

\begin{commande}{exe, exo, prb, rem}{carnomath}

\Description{
Commandes pour exemples, exercices, problèmes et remarques.
}

\Syntaxe{
\syntaxe{\exe{contenu}}
\syntaxe{\exo{contenu}}
\syntaxe{\prb{contenu}}
\syntaxe{\rem{contenu}}
}

\Arguments{
Un argument obligatoire.
}

\Exemple{
\exe{
La fonction $f(x)=2x+3$ est une fonction affine.
}

\exo{
Résoudre l'équation $2x+3=7$.
}

\rem{
Le coefficient directeur peut être nul.
}
}

\end{commande}

\begin{commande}{dem, sol, cqfd}{carnomath}

\Description{
Commandes pour démonstrations, solutions et fin de preuve.
}

\Syntaxe{
\syntaxe{\dem{contenu}}
\syntaxe{\sol{contenu}}
\syntaxe{\cqfd}
}

\Arguments{
\cmd{dem} et \cmd{sol} prennent un argument.

\cmd{cqfd} ne prend aucun argument.
}

\Exemple{
\dem{
On développe puis on simplifie.

\cqfd
}

\sol{
On obtient finalement $x=2$.
}
}

\end{commande}

%%%%%%%%%%%%%%%%%%%%%%%%%%%%%%%%%%%%%%%%%%%%%%%%%%%%%%%%%%%%%%%%%%%%%%%%%%%%%%
\section{Pédagogie}
%%%%%%%%%%%%%%%%%%%%%%%%%%%%%%%%%%%%%%%%%%%%%%%%%%%%%%%%%%%%%%%%%%%%%%%%%%%%%%

\begin{commande}{level}{carnopeda}

\Description{
Sélectionne la version du document.
}

\Syntaxe{
\syntaxe{\level{eleve}}
\syntaxe{\level{prof}}
\syntaxe{\level{notes}}
}

\Arguments{
Un argument obligatoire : \texttt{eleve}, \texttt{prof} ou \texttt{notes}.
}

\end{commande}

\begin{commande}{obj, pre, idee, oral, tab, quest, piege, retenir}{carnopeda}

\Description{
Panneaux pédagogiques principaux.
}

\Syntaxe{
\syntaxe{\obj{contenu}}
\syntaxe{\pre{contenu}}
\syntaxe{\idee{contenu}}
\syntaxe{\oral{contenu}}
\syntaxe{\tab{contenu}}
\syntaxe{\quest{contenu}}
\syntaxe{\piege{contenu}}
\syntaxe{\retenir{contenu}}
}

\Arguments{
Un argument obligatoire contenant le texte du panneau.
}

\Exemple{
\obj{
Comprendre la notion de spectre.
}

\idee{
Relier série de Fourier et transformée de Fourier.
}

\retenir{
La transformée de Fourier analyse les fréquences d'un signal.
}
}

\end{commande}

\begin{commande}{note, todo, source, experience, temps}{carnopeda}

\Description{
Commandes de préparation et d'annotation pour l'enseignant.
}

\Syntaxe{
\syntaxe{\note{contenu}}
\syntaxe{\note[titre]{contenu}}
\syntaxe{\todo{contenu}}
\syntaxe{\source[titre]{contenu}}
\syntaxe{\experience[titre]{contenu}}
\syntaxe{\temps{minutes}}
}

\Arguments{
Selon les commandes, un titre optionnel et un contenu obligatoire.

\cmd{temps} prend une durée en minutes.
}

\Exemple{
\note[Préparation]{
Prévoir une animation GeoGebra.
}

\todo[Figure]{
Créer une figure de spectre.
}

\temps{10}
}

\end{commande}

\begin{commande}{exobandeau}{carnopeda}

\Description{
Crée un bandeau d'exercice numéroté, avec difficulté et points éventuels en marge.
}

\Syntaxe{
\syntaxe{\exobandeau{titre}{difficulté}}
\syntaxe{\exobandeau[points]{titre}{difficulté}}
}

\Arguments{
\begin{itemize}
  \item \texttt{points} : nombre de points optionnel ;
  \item \texttt{titre} : titre de l'exercice ;
  \item \texttt{difficulté} : entier de 1 à 5.
\end{itemize}
}

\Exemple{
\exobandeau{Application directe}{2}

Énoncé de l'exercice.

\exobandeau[5]{Problème guidé}{4}

Énoncé de l'exercice.
}

\end{commande}

%%%%%%%%%%%%%%%%%%%%%%%%%%%%%%%%%%%%%%%%%%%%%%%%%%%%%%%%%%%%%%%%%%%%%%%%%%%%%%
\section{Compatibilité}
%%%%%%%%%%%%%%%%%%%%%%%%%%%%%%%%%%%%%%%%%%%%%%%%%%%%%%%%%%%%%%%%%%%%%%%%%%%%%%

\cadre{CarnoWarning}[carnolegacy]{
Le package \pkg{carnolegacy} est destiné uniquement à compiler les anciens documents.

Les commandes documentées ici ne sont pas recommandées pour les nouveaux documents.
}

\begin{commande}{Cbox}{carnolegacy}

\Description{
Ancienne commande de bloc coloré conservée pour compatibilité.
}

\Syntaxe{
\syntaxe{\Cbox{couleur}{titre}{contenu}}
}

\Remarques{
Pour les nouveaux documents, préférer \cmd{cadre}.
}

\end{commande}

\begin{commande}{red}{carnolegacy}

\Description{
Ancienne commande de mise en couleur rouge.
}

\Syntaxe{
\syntaxe{\red{texte}}
}

\Arguments{
Un argument contenant le texte à afficher en rouge.
}

\Exemple{
\red{Texte mis en rouge.}
}

\Remarques{
Cette commande est conservée uniquement pour la compatibilité avec les anciens documents.
}

\end{commande}

%%%%%%%%%%%%%%%%%%%%%%%%%%%%%%%%%%%%%%%%%%%%%%%%%%%%%%%%%%%%%%%%%%%%%%%%%%%%%%
\section{Annexe : tableau de l'API}
%%%%%%%%%%%%%%%%%%%%%%%%%%%%%%%%%%%%%%%%%%%%%%%%%%%%%%%%%%%%%%%%%%%%%%%%%%%%%%

\begin{center}
\renewcommand{\arraystretch}{1.3}
\begin{tabular}{|l|l|c|c|}
\hline
\textbf{Commande} & \textbf{Package} & \textbf{Testée} & \textbf{Documentée} \\
\hline
\cmd{cadre} & carnotex & oui & oui \\
\cmd{col} & carnotex & oui & oui \\
\cmd{bordure} & carnotex & oui & oui \\
\cmd{colorline} & carnotex & oui & oui \\
\cmd{image} & carnotex & oui & oui \\
\cmd{df}, \cmd{prop}, \cmd{thrm} & carnomath & oui & oui \\
\cmd{obj}, \cmd{note}, \cmd{todo} & carnopeda & oui & oui \\
\cmd{Cbox} & carnolegacy & oui & oui \\
\hline
\end{tabular}
\end{center}

\end{document}